
With over 7,635 benchmarks across 4,886 LLM evaluation papers showing less than 25\% overlap, researchers face a benchmark selection problem: how can they prospectively identify reliable benchmarks before committing experimental resources? This question has practical consequences. Without prospective verification tools, researchers risk investing in benchmarks that produce unstable cross-benchmark rankings.

This study tests whether coefficient of variation (CV), a zero-cost statistic computable from any published leaderboard, predicts cross-benchmark ranking stability. The hypothesis is that high CV indicates inconsistent model differentiation, suggesting measurement noise that should reduce stability when rankings are compared across benchmarks. If validated, CV would provide a prospective quality signal for benchmark selection.

Using a mock benchmark corpus of 10 trust evaluation benchmarks (15 models each), we find CV shows weak, non-significant correlation with mean cross-benchmark Spearman $\rho$ (Pearson r=-0.486, p=0.154, 95\% CI: [-0.854, 0.207]). This null result fails pre-registered criteria (r<-0.5, p<0.05) by narrow margins---the correlation magnitude is 6\% short of threshold, and the p-value is $3\times$ above the significance level. The direction is negative as hypothesized, but the effect is too weak and uncertain to serve as a reliable quality predictor.

However, cross-benchmark correlation analysis reveals negative and near-zero correlations between benchmarks ostensibly measuring ''trustworthiness.'' FaithfulQA and FinTrust show $\rho$=-0.568, while many pairs exhibit near-zero correlations (e.g., FinTrust-HaluBench $\rho$=-0.007). These patterns suggest trust benchmarks may measure orthogonal sub-dimensions (hallucination detection, financial ethics, factual accuracy) rather than a unitary construct, confounding what ''cross-benchmark stability'' means as a quality metric.

This work contributes: (1) rigorous negative evidence that simple variance metrics are insufficient for benchmark verification, (2) documentation of cross-benchmark correlation patterns revealing construct validity issues, and (3) a methodological framework for testing benchmark quality predictors with pre-registered falsification criteria.

A critical limitation is the use of mock benchmark data. All values are synthetic, not extracted from real TrustLLM, HaluBench, or TruthfulQA leaderboards. Real-leaderboard replication is required to validate findings before external validity can be claimed.
